\section{Introduction}
\label{sec:intro}

\subsection{Motivation}
Documents on the web fall broadly into two categories. \emph{Text-based} PDFs
embed their characters as selectable, searchable glyphs and work seamlessly with
the browser's built-in find feature. \emph{Image-based} documents---scanned
pages, photographs of text, screenshots, and PDFs whose text has been flattened
into raster images---contain no machine-readable text at all. To the browser
they are pictures. Pressing $\mathrm{Ctrl}{+}\mathrm{F}$ on such a document finds
nothing, even though the text is plainly visible to a human reader.

This is a pervasive friction point. A student searching a scanned textbook, a
researcher scanning archival material, or anyone working with a faxed or
photographed form repeatedly hits the same wall: the content is \emph{there}, but
not \emph{findable}. The conventional remedy---uploading the file to a
server-side OCR service to produce a new, searchable PDF---is slow, raises
privacy concerns, and breaks the natural in-browser reading flow.

\subsection{Key Idea}
\proj{} takes a deliberately different stance. Instead of \emph{reinventing}
search---building an index, a query parser, and highlight logic---it
\emph{transforms the document so the browser can search it natively}. Every page
is rendered, any text baked into images is recovered with OCR, and the recovered
text is injected back into the page as invisible, exactly-positioned DOM
elements. From the browser's perspective the document now \emph{is} text, and the
entire mature machinery of in-page search applies for free.

This inversion is the conceptual core of the project: \emph{the cheapest way to
make something searchable is to make it real text}, not to intercept and
re-implement the act of searching.

\subsection{Contributions}
This report documents the design and implementation of \proj{}. Its principal
contributions are:

\begin{itemize}
  \item A fully \textbf{client-side, offline} pipeline that makes scanned PDFs
        searchable in the browser, with no server round-trip and no data leaving
        the user's machine.
  \item A \textbf{coordinate-mapping engine} that aligns OCR output and native
        PDF text across four distinct coordinate spaces, remaining correct under
        zoom, page rotation, and high-DPI rendering
        (Section~\ref{sec:design}).
  \item A \textbf{performance architecture}---lazy, viewport-driven rendering
        plus idle-time background indexing on Web Workers, with
        resolution-independent caching in IndexedDB---that keeps the interface
        responsive on large documents (Section~\ref{sec:impl}).
  \item An \textbf{OCR accuracy layer}: high-resolution rendering with targeted
        image preprocessing, plus post-processing that emits search-only
        correction variants to raise the find hit-rate
        (Section~\ref{sec:impl}).
\end{itemize}

\subsection{Report Structure}
Section~\ref{sec:requirements} states the system's functional requirements as
user stories. Section~\ref{sec:background} surveys the relevant background and
prior approaches. Section~\ref{sec:tech} catalogues the technologies used and how
they coordinate. Section~\ref{sec:design} presents the system architecture and the
coordinate-mapping problem at its heart. Section~\ref{sec:process} recounts the
phased development of the project and the reasoning behind each phase, and
Section~\ref{sec:impl} details the implementation of each pipeline stage.
Section~\ref{sec:results} evaluates the system qualitatively,
Section~\ref{sec:limitations} discusses limitations and future work, and
Section~\ref{sec:conclusion} concludes.
